\subsection{A complete presentation by adjacent crossings}

\begin{result}{E.38}{Proved}{Endpoint actions and transported smoothing diamonds give every relation}
The discrete exit category is presented by its endpoint automorphism
groups, its adjacent crossing bisets with their actual stabilizer
maps, and the smoothing diamonds obtained from actual pinned
multicurves by deleting two distinct curves in either order.
All pins and intermediate markings are transported as specified below.
There is no further relation at the level of the ordinary category.

Its strict finite-congruence completion is the universal profinite
category on the completed endpoint groups and adjacent bisets with
those same continuously transported diamond equations.  This is a
presentation in terms of its finite categorical quotients.  The
comparison and inverse are induced by the displayed normal forms
and agree on every endpoint, adjacent arrow and specified flag.
In particular this result uses no assertion about the normalizer of
an arbitrary completed multicurve stabilizer.
\end{result}
\begin{proof}
We first work before taking an orbit skeleton.  Fix the surface and
write $\Gamma$ for its pure orientation-preserving mapping class
group.  Let $\mathcal C$ have all actual isotopy classes of
multicurves as objects, including the empty multicurve, and
\[
 \Hom_{\mathcal C}(\sigma,\tau)
 =\{g\in\Gamma:g\sigma\supseteq\tau\}/\Delta_\sigma,
 \qquad \Delta_\sigma=\langle T_c:c\in\sigma\rangle.
                                                               \tag{E.140}
\]
The quotient is on the right, and composition is multiplication:
$[h]\circ[g]=[hg]$.  The twist groups are generated by positive
twists.  An element of $\Delta_\sigma$ fixes every curve of
$\sigma$, so the displayed quotient and composition are defined
as in A.4.

Use the following generators, with their full typed domains:
\[
 u_{g,\sigma}:\sigma\longrightarrow g\sigma
       \quad(g\in\Gamma),\qquad
 d_{\sigma,c}:\sigma\longrightarrow\sigma\setminus\{c\}
       \quad(c\in\sigma).                                    \tag{E.141}
\]
The $u$ are invertible.  Impose precisely the following relations:
\[
\begin{aligned}
 u_{h,g\sigma}u_{g,\sigma}&=u_{hg,\sigma},
       &u_{1,\sigma}&=1,\\
 u_{\delta,\sigma}&=1\quad(\delta\in\Delta_\sigma),\\
 u_{g,\sigma\setminus\{c\}}d_{\sigma,c}
       &=d_{g\sigma,gc}u_{g,\sigma},\\
 d_{\sigma\setminus\{c\},d}\,d_{\sigma,c}
       &=d_{\sigma\setminus\{d\},c}\,d_{\sigma,d}
       \quad(c\ne d\text{ in }\sigma).
\end{aligned}                                                   \tag{E.142}
\]
Products in (E.142) use categorical order, with the rightmost
arrow traversed first.  The evident functor to (E.140) sends
$u_{g,\sigma}$ to $[g]$ and $d_{\sigma,c}$ to $[1]$.
Every displayed relation holds there.

Here is its inverse normal form, including the equality test.
For $\rho\subseteq\sigma$, let $D_{\sigma,\rho}$ be the
successive deletion of the curves of $\sigma\setminus\rho$.
Any two orders differ by adjacent transpositions, and the last
relation of (E.142) is exactly each such transposition.  Thus
$D_{\sigma,\rho}$ is well-defined.  Slide every $u$ to the end
of a word by the equivariance relation, then combine the $u$'s.
Every resulting arrow has the form
\[
                u_{g,\rho}D_{\sigma,\rho},
       \qquad \rho\subseteq\sigma,\quad g\rho=\tau.           \tag{E.143}
\]
Conversely, for $[g]\in\Hom(\sigma,\tau)$ take
$\rho=g^{-1}\tau$ to obtain (E.143).

Suppose the two resulting transporter classes are equal.
Then $h=g\delta$ for some $\delta\in\Delta_\sigma$.
Since $g^{-1}\tau$ is a subset of $\sigma$, the twist $\delta$
fixes it, so both forms have the same $\rho$.  Iterated
equivariance and the source twist relation give
\[
 \begin{aligned}
 u_{h,\rho}D_{\sigma,\rho}
 &=u_{g,\rho}u_{\delta,\rho}D_{\sigma,\rho}\\
 &=u_{g,\rho}D_{\sigma,\rho}u_{\delta,\sigma}
  =u_{g,\rho}D_{\sigma,\rho}.
 \end{aligned}                                                 \tag{E.144}
\]
This is precisely the right quotient by $\Delta_\sigma$,
including a twist about a curve that has been deleted.
Conversely (E.142) cannot change the transporter coset,
because each relation holds in (E.140).  Hence these are
mutually inverse functors, and the presentation is complete.
This proves the relation statement directly, without a
normalizer or a connectivity assumption on a curve complex.

We now descend it to finite orbit data.  Choose one actual
representative $\sigma_i$ of each multicurve type.  For every
multicurve $\sigma$ choose $k_\sigma\in\Gamma$ with
$k_\sigma\sigma=\sigma_i$, taking $k_{\sigma_i}=1$.
The reduction functor and its section are explicit:
\[
 R(\sigma)=i,\qquad
 R([g]:\sigma\to\tau)
       =[k_\tau g k_\sigma^{-1}]:i\to j.                     \tag{E.145}
\]
Conjugation carries $\Delta_\sigma$ to $\Delta_{\sigma_i}$,
so (E.145) is well-defined on the right cosets.  Its
composition law follows by cancellation of the middle
$k_\tau^{-1}k_\tau$.  Inclusion $J$ of the representatives
satisfies $RJ=1$; the arrows $u_{k_\sigma^{-1},\sigma_i}$
give the natural isomorphism $JR\Rightarrow1$.

There are finitely many types, adjacent incidence types,
and ordered two-curve deletion types: they are finite stable
labelled graphs with a specified retained edge subset or a
specified ordered pair of edges.  Retain every incidence
occurrence, including multiple occurrences with the same
source and target graph type.  For a chosen actual adjacent
flag $\sigma\supset\tau$, its crossing stabilizer and endpoint
maps are
\[
 K=\frac{\Stab(\sigma)\cap\Stab(\tau)}{\Delta_\tau}
 \longrightarrow
 \begin{cases}
 G_1=\Stab(\tau)/\Delta_\tau,\quad\text{inclusion},\\
 G_0=\Stab(\sigma)/\Delta_\sigma,\quad\text{quotient}.
 \end{cases}                                                   \tag{E.146}
\]
Its adjacent biset orbit is the transitive
$(G_1\times G_0)$-set with that stabilizer, where the action
is $a\mapsto g_1ag_0^{-1}$.  For a representative whose
retained multicurve is transported by $k_\tau$, conjugate
the target map in (E.146) by $k_\tau$ as prescribed in
(E.145).  The source map is unchanged.  Thus the data are
exactly the adjacent bisets and maps of A.6 and E.32.
Their stabilizer equations express precisely the changes
of endpoint trivialization; no full stabilizer is substituted
for an incidence stabilizer.

The matching rule for every diamond is equally explicit.
Choose an actual multicurve $\sigma$ and two distinct actual
curves $c,d\in\sigma$.  Use the four vertices
$\sigma$, $\sigma\setminus c$, $\sigma\setminus d$,
$\rho=\sigma\setminus\{c,d\}$, with their chosen
transports $k$ in (E.145).  The four skeletal adjacent
arrows are the classes of
\[
 \begin{array}{ll}
 k_{\sigma\setminus c}k_\sigma^{-1},&
 k_\rho k_{\sigma\setminus c}^{-1},\\
 k_{\sigma\setminus d}k_\sigma^{-1},&
 k_\rho k_{\sigma\setminus d}^{-1}.
 \end{array}
\]
Their first composite is
\[
 (k_\rho k_{\sigma\setminus c}^{-1})
 (k_{\sigma\setminus c}k_\sigma^{-1})
   =k_\rho k_\sigma^{-1}.                                    \tag{E.147}
\]
The other order has the identical product.  Express each
of the four classes in its chosen adjacent biset using
these actual endpoint transports, and impose their equality.
Changing a representative of an endpoint orbit changes the
expression by exactly (E.146).  In particular the retained
pins in the two routes are transported from this same
actual $\sigma,c,d$; matching only the abstract endpoint
graph classes would omit part of (E.147).

A longer flag is treated by transporting every vertex of
one actual nested chain through its own $k_\sigma$.
Its composites telescope by (E.145).  Any two orders of
its deletions are related by the adjacent exchanges just
proved.  Hence its equations follow from the transported
diamonds and endpoint actions.  Conversely every skeletal
equation lifts, using $J$ and the displayed natural
isomorphism, to the complete presentation (E.142).
This proves both completeness and independence of all
chosen representatives at the ordinary-category level.
It is not an assertion that higher homotopy coherence is
lost when considering higher-valued coefficient systems.

It remains to prove the completion statement with its
actual finite congruences.  Complete every endpoint group
and every adjacent biset using E.32.  For each of the
finitely many actual diamond representatives (E.147),
take all its endpoint translates and their closure in
the product of its four completed adjacent Hom-sets.
Equivalently apply the product of its four completed
endpoint groups to the stated quadruple.  That compact
orbit is the closure of the discrete orbit, because
the discrete endpoint groups are dense.  Take the
finite union of these transport families.  This specifies
which completed flags are matched; it imposes no equation
on a pair of unrelated flags merely because their
intermediate graph classes coincide.

Consider finite categorical assignments of the following
data: finite quotients of the endpoint groups, continuous
maps from the adjacent bisets to the corresponding finite
Hom-sets, equivariance for both endpoint actions, and the
just specified diamond equalities.  Such an assignment
restricts to the discrete generators.  The normal-form
proof (E.142)--(E.144), followed by (E.145), extends it
uniquely to a functor on the discrete skeletal category.
Its kernel on the Hom-sets is a finite categorical
congruence.  Conversely every finite functor on that
category extends uniquely to the completed endpoints
and adjacent bisets by E.32.  Its diamond equations
hold on the dense discrete transport families and hence
on their closures, since equality in a finite set is
closed.  After replacing each finite target by its image and identifying
isomorphic image targets, these operations are inverse, including
refinement maps between finite quotients.  In symbols,
\[
 \left\{\begin{gathered}
 \text{compatible finite}\\
 \text{endpoint/biset assignments}
 \end{gathered}\right\}
 \simeq
 \left\{\begin{gathered}
 \text{finite categorical congruence}\\
 \text{quotients of }C
 \end{gathered}\right\}.
                                                               \tag{E.148}
\]
One may replace a finite target category by the image of
the extended functor; this keeps the object set fixed
and gives a quotient in exactly the sense of E.1.
Intersections give common refinements on both sides.
Thus (E.148) also proves the required cofinality.

Define the universal profinite presentation as the inverse
limit of these finite assignments.  The same inverse
system is E.1's strict completion, so its comparison is
\[
 \mathcal P_{\mathrm{adj},\diamond}
                  \xrightarrow{\ \sim\ }\widehat C.           \tag{E.149}
\]
On generators the map is the actual endpoint or adjacent
inclusion.  Its inverse on every finite coordinate is
the normal-form extension in (E.148); these coordinates
are compatible with refinement and therefore define
the inverse continuous functor.  Their composites fix
every finite coordinate and hence every morphism.
All arrows are covered: a discrete arrow is a composite
of adjacent arrows and endpoint isomorphisms by (E.143);
for each noninvertible completed Hom-set, the finite
union of the images of the corresponding compact
products is closed and contains its dense discrete
Hom-set.  Endpoint automorphisms are already among
the generators.  This also describes (E.149) by actual
completed factorizations.

The proof uses the universal finite-quotient presentation
just defined.  It does not replace it by an unjustified
ordinary quotient of a compact arrow space by a closure
of relations.  Equations (E.148)--(E.149) give the exact
comparison required for the finite-congruence completion,
with all full-flag transport data retained.
\end{proof}
